\documentclass[12pt,a4paper]{article}

\usepackage[utf8]{inputenc}
\usepackage[T1]{fontenc}
\usepackage[serbian]{babel}
\usepackage{geometry}
\usepackage{amsmath, amssymb}
\usepackage{booktabs}
\usepackage{longtable}
\usepackage{array}
\usepackage{hyperref}
\usepackage{listings}
\usepackage{xcolor}
\usepackage{tcolorbox}
\usepackage{xcolor} % za boje
\geometry{margin=2.5cm}

\hypersetup{
    colorlinks=true,
    linkcolor=blue,
    urlcolor=blue,
    citecolor=blue
}

\lstset{
  basicstyle=\ttfamily\small,
  keywordstyle=\color{blue},
  commentstyle=\color{gray},
  stringstyle=\color{teal},
  numbers=left,
  numberstyle=\tiny\color{gray},
  stepnumber=1,
  numbersep=8pt,
  breaklines=true,
  frame=single,
  tabsize=2
}


% Definicija crvenog box-a za napomene
\tcbset{
    myrednote/.style={
        colback=red!5,        % svetlo crvena pozadina
        colframe=red!75!black, % tamno crvena ivica
        fonttitle=\bfseries,
        coltitle=black!80!black,
        boxrule=0.8pt,
        arc=4pt,
        left=2mm, right=2mm, top=1mm, bottom=1mm
    }
}


\title{\textbf{TS-FCTP (Two-Stage Fixed Charge Transportation Problem)}\\
Seminarski rad -- Tema broj 8}
\author{
Mateja Stojanović (mi251107)
}
\date{}

\begin{document}
\maketitle

\section*{Osnovne informacije}
\begin{itemize}
    \item \textbf{Ime i prezime:} Mateja Stojanović
    \item \textbf{Broj indeksa:} mi251107
    \item \textbf{Tema/Zadatak:} Tema broj 8
\end{itemize}

\section{Opis problema}
Problem TS-FCTP (Two-Stage Fixed Charge Transportation Problem) predstavlja problem transporta u dve faze kroz mrežu sa tri nivoa čvorova:
\begin{itemize}
    \item proizvođači (izvori) $i \in I$,
    \item distributivni centri $j \in J$,
    \item kupci $k \in K$.
\end{itemize}

Roba se transportuje u dve etape:
\[
i \rightarrow j \rightarrow k.
\]

Za svaku vezu u mreži postoje:
\begin{itemize}
    \item \textbf{varijabilni trošak} po jedinici protoka,
    \item \textbf{fiksni trošak} koji se plaća ukoliko se veza koristi (tj. ukoliko se šalje neka količina robe).
\end{itemize}

Cilj je minimizacija ukupnog troška transporta uz poštovanje ograničenja kapaciteta izvora i ispunjenje potražnje kupaca.

\section{Matematička formulacija problema}

\subsection{Skupovi i parametri}
\begin{itemize}
    \item $I$ -- skup proizvođača (izvora), $|I| = I$,
    \item $J$ -- skup distributivnih centara, $|J| = J$,
    \item $K$ -- skup kupaca, $|K| = K$,
    \item $s_i$ -- raspoloživa ponuda izvora $i$,
    \item $d_k$ -- potražnja kupca $k$,
    \item $b_{ij}$ -- varijabilni trošak transporta od izvora $i$ do centra $j$,
    \item $f_{ij}$ -- fiksni trošak aktiviranja veze $(i,j)$,
    \item $c_{jk}$ -- varijabilni trošak transporta od centra $j$ do kupca $k$,
    \item $g_{jk}$ -- fiksni trošak aktiviranja veze $(j,k)$.
\end{itemize}

\subsection{Promenljive}
\begin{itemize}
    \item $x_{ij} \ge 0$ -- količina koja se transportuje od izvora $i$ do centra $j$,
    \item $y_{jk} \ge 0$ -- količina koja se transportuje od centra $j$ do kupca $k$,
    \item $z_{ij} \in \{0,1\}$ -- indikator da li se koristi veza $(i,j)$,
    \item $w_{jk} \in \{0,1\}$ -- indikator da li se koristi veza $(j,k)$.
\end{itemize}

\subsection{Funkcija cilja}
Minimizuje se suma varijabilnih i fiksnih troškova u obe faze:
\[
\min \sum_{i \in I}\sum_{j \in J} \left( b_{ij}x_{ij} + f_{ij}z_{ij} \right)
+ \sum_{j \in J}\sum_{k \in K} \left( c_{jk}y_{jk} + g_{jk}w_{jk} \right).
\]

\subsection{Ograničenja}
\paragraph{Ograničenje ponude izvora}
\[
\sum_{j \in J} x_{ij} \le s_i, \quad \forall i \in I.
\]

\paragraph{Zadovoljenje potražnje kupaca}
\[
\sum_{j \in J} y_{jk} = d_k, \quad \forall k \in K.
\]

\paragraph{Balans protoka u distributivnim centrima}
Ukupni ulaz u centar mora biti jednak ukupnom izlazu:
\[
\sum_{i \in I} x_{ij} - \sum_{k \in K} y_{jk} = 0, \quad \forall j \in J.
\]

\paragraph{Aktivacija fiksnih troškova (Big-M povezivanje)}
Ako se veza koristi, binarna promenljiva mora biti 1:
\[
x_{ij} \le s_i z_{ij}, \quad \forall i \in I, \forall j \in J,
\]
\[
y_{jk} \le d_k w_{jk}, \quad \forall j \in J, \forall k \in K.
\]

\section{Značenje funkcije cilja i uslova}
\begin{itemize}
    \item \textbf{Funkcija cilja} predstavlja ukupni trošak transporta:
    \begin{itemize}
        \item $b_{ij}x_{ij}$ i $c_{jk}y_{jk}$ su varijabilni troškovi proporcionalni količini robe,
        \item $f_{ij}z_{ij}$ i $g_{jk}w_{jk}$ su fiksni troškovi koji se plaćaju jednom ako se veza koristi.
    \end{itemize}
    \item \textbf{Ograničenja ponude} obezbeđuju da nijedan izvor ne šalje više od raspoložive količine.
    \item \textbf{Ograničenja potražnje} garantuju da se zahtevi svih kupaca tačno zadovolje.
    \item \textbf{Balans u centrima} znači da distributivni centri ne mogu praviti zalihe niti stvarati robu, već samo prosleđuju protok.
    \item \textbf{Big-M ograničenja} povezuju kontinualne tokove sa binarnim promenljivama, čime se modeluje fiksni trošak aktivacije veze.
\end{itemize}

\section{Instance korišćene u eksperimentu}

\subsection{Podela instanci (male, srednje, velike, baseline)}
Instance su generisane u tri grupe na osnovu dimenzija $(I,J,K)$:

\subsubsection*{Male instance (20 instanci)}
\begin{center}
\begin{tabular}{lll}
\toprule
\textbf{Instanca} & \textbf{Dimenzije} $(I,J,K)$ \\
\midrule
small\_1  & (4,8,16) \\
small\_2  & (4,8,20) \\
small\_3  & (4,10,20) \\
small\_4  & (6,12,24) \\
small\_5  & (6,12,30) \\
small\_6  & (6,15,30) \\
small\_7  & (8,16,32) \\
small\_8  & (8,16,40) \\
small\_9  & (8,20,40) \\
small\_10 & (10,20,40) \\
small\_11 & (10,20,50) \\
small\_12 & (10,25,50) \\
small\_13 & (12,24,48) \\
small\_14 & (12,24,60) \\
small\_15 & (12,30,60) \\
small\_16 & (14,28,56) \\
small\_17 & (14,28,70) \\
small\_18 & (16,32,64) \\
small\_19 & (16,32,80) \\
small\_20 & (18,36,72) \\
\bottomrule
\end{tabular}
\end{center}

\subsubsection*{Srednje instance (20 instanci)}
\begin{center}
\begin{tabular}{lll}
\toprule
\textbf{Instanca} & \textbf{Dimenzije} $(I,J,K)$ \\
\midrule
medium\_1  & (20,40,80) \\
medium\_2  & (20,40,100) \\
medium\_3  & (20,50,100) \\
medium\_4  & (25,50,100) \\
medium\_5  & (25,50,125) \\
medium\_6  & (25,60,120) \\
medium\_7  & (30,60,120) \\
medium\_8  & (30,60,150) \\
medium\_9  & (30,75,150) \\
medium\_10 & (35,70,140) \\
medium\_11 & (35,70,175) \\
medium\_12 & (40,80,160) \\
medium\_13 & (40,80,200) \\
medium\_14 & (45,90,180) \\
medium\_15 & (45,90,225) \\
medium\_16 & (50,100,200) \\
medium\_17 & (50,100,250) \\
medium\_18 & (55,110,220) \\
medium\_19 & (55,110,275) \\
medium\_20 & (60,120,240) \\
\bottomrule
\end{tabular}
\end{center}

\subsubsection*{Velike instance (20 instanci)}
\begin{center}
\begin{tabular}{lll}
\toprule
\textbf{Instanca} & \textbf{Dimenzije} $(I,J,K)$ \\
\midrule
large\_1  & (70,140,280) \\
large\_2  & (70,140,350) \\
large\_3  & (80,160,320) \\
large\_4  & (80,160,400) \\
large\_5  & (90,180,360) \\
large\_6  & (90,180,450) \\
large\_7  & (100,200,400) \\
large\_8  & (100,200,500) \\
large\_9  & (110,220,440) \\
large\_10 & (110,220,550) \\
large\_11 & (120,240,480) \\
large\_12 & (120,240,600) \\
large\_13 & (130,260,520) \\
large\_14 & (130,260,650) \\
large\_15 & (140,280,560) \\
large\_16 & (140,280,700) \\
large\_17 & (150,300,600) \\
large\_18 & (150,300,750) \\
large\_19 & (160,320,640) \\
large\_20 & (160,320,800) \\
\bottomrule
\end{tabular}
\end{center}

\subsection{Baseline instance}
Pored tri osnovne grupe instanci (male, srednje i velike), generisana je i posebna grupa \textit{baseline} instanci.
Ove instance služe kao referentni skup za poređenje performansi solvera, jer imaju umerene dimenzije.

Baseline grupa sadrži 4 instance sledećih dimenzija $(I,J,K)$:

\begin{center}
\begin{tabular}{ll}
\toprule
\textbf{Instanca} & \textbf{Dimenzije} $(I,J,K)$ \\
\midrule
baseline\_1 & (10,20,40) \\
baseline\_2 & (15,25,50) \\
baseline\_3 & (20,30,60) \\
baseline\_4 & (25,40,80) \\
\bottomrule
\end{tabular}
\end{center}


\subsection{Primer jedne manje instance}
Format instance fajla je:
\begin{itemize}
    \item prva linija: $I\ J\ K$
    \item druga linija: ponude $s_i$
    \item treća linija: potražnje $d_k$
    \item zatim matrice troškova: $b_{ij}$, $f_{ij}$, $c_{jk}$, $g_{jk}$
\end{itemize}

Primer (struktura) male instance \texttt{small\_1.txt} dimenzija $(4,8,16)$:
\begin{verbatim}
4 8 16
80 80 80 80
... (16 vrednosti potraznje koje zbirno daju 320)
... (matrica b_ij dim 4x8)
... (matrica f_ij dim 4x8)
... (matrica c_jk dim 8x16)
... (matrica g_jk dim 8x16)
\end{verbatim}

\section{Generisanje instanci}
Instance su generisane programom \texttt{generator.cpp}. Koordinate čvorova se nasumično biraju iz opsega $[-400,400]$,
a rastojanja se računaju Euklidskom metrikom i zaokružuju na dole.

Ponude su fiksne:
\[
s_i = 80 \quad \forall i \in I,
\]
dok su potražnje inicijalno generisane uniformno u opsegu $[10,30]$, a zatim se koriguju tako da važi:
\[
\sum_{k \in K} d_k = \sum_{i \in I} s_i.
\]

Varijabilni troškovi su proporcionalni rastojanju, dok se fiksni troškovi dobijaju množenjem rastojanja slučajnim faktorom:
\begin{itemize}
    \item za prvu fazu faktor $r_f \in [10,15]$,
    \item za drugu fazu faktor $r_g \in [5,10]$.
\end{itemize}

\subsection*{Kod generatora instanci (\texttt{generator.cpp})}
\begin{lstlisting}[language=C++]
#include <fstream>
#include <vector>
#include <cmath>
#include <random>
#include <string>

struct Node {
    int x, y;
};

int dist(const Node& a, const Node& b) {
    return (int)std::floor(
        std::sqrt((a.x - b.x)*(a.x - b.x) +
                  (a.y - b.y)*(a.y - b.y))
        );
}

void generate_group(const std::string& prefix,
                    const std::vector<std::vector<int>>& dims,
                    std::mt19937& rng) {

    std::uniform_int_distribution<int> coord(-400,400);
    std::uniform_int_distribution<int> rf(10,15);
    std::uniform_int_distribution<int> rg(5,10);
    std::uniform_int_distribution<int> dem(10,30);

    for (int t = 0; t < (int)dims.size(); t++) {
        int I = dims[t][0];
        int J = dims[t][1];
        int K = dims[t][2];

        std::vector<Node> P(I), D(J), C(K);
        for (auto& x : P) x = {coord(rng), coord(rng)};
        for (auto& x : D) x = {coord(rng), coord(rng)};
        for (auto& x : C) x = {coord(rng), coord(rng)};

        std::vector<int> s(I, 80);
        std::vector<int> d(K);
        int supply = I * 80, demand = 0;

        for (int k = 0; k < K; k++) {
            d[k] = dem(rng);
            demand += d[k];
        }

        for (int k = 0; demand != supply; k = (k+1)%K) {
            if (demand > supply && d[k] > 1) {
                d[k]--; demand--;
            } else if (demand < supply) {
                d[k]++; demand++;
            }
        }

        std::ofstream out("generated_instances/" + prefix +
                          "_" + std::to_string(t+1) + ".txt");

        out << I << " " << J << " " << K << "\n";

        for (int i = 0; i < I; i++) out << s[i] << " ";
        out << "\n";

        for (int k = 0; k < K; k++) out << d[k] << " ";
        out << "\n";

        for (int i = 0; i < I; i++) {
            for (int j = 0; j < J; j++)
                out << dist(P[i], D[j]) << " ";
            out << "\n";
        }

        for (int i = 0; i < I; i++) {
            int r = rf(rng);
            for (int j = 0; j < J; j++)
                out << dist(P[i], D[j]) * r << " ";
            out << "\n";
        }

        for (int j = 0; j < J; j++) {
            for (int k = 0; k < K; k++)
                out << dist(D[j], C[k]) << " ";
            out << "\n";
        }

        for (int j = 0; j < J; j++) {
            int r = rg(rng);
            for (int k = 0; k < K; k++)
                out << dist(D[j], C[k]) * r << " ";
            out << "\n";
        }

        out.close();
    }
}

int main() {

    std::mt19937 rng(1389);

    std::vector<std::vector<int>> small = {
        {4,8,16}, {4,8,20}, {4,10,20}, {6,12,24}, {6,12,30},
        {6,15,30}, {8,16,32}, {8,16,40}, {8,20,40}, {10,20,40},
        {10,20,50}, {10,25,50}, {12,24,48}, {12,24,60}, {12,30,60},
        {14,28,56}, {14,28,70}, {16,32,64}, {16,32,80}, {18,36,72}
    };

    std::vector<std::vector<int>> medium = {
        {20,40,80}, {20,40,100}, {20,50,100}, {25,50,100}, {25,50,125},
        {25,60,120}, {30,60,120}, {30,60,150}, {30,75,150}, {35,70,140},
        {35,70,175}, {40,80,160}, {40,80,200}, {45,90,180}, {45,90,225},
        {50,100,200}, {50,100,250}, {55,110,220}, {55,110,275}, {60,120,240}
    };

    std::vector<std::vector<int>> large = {
        {70,140,280}, {70,140,350}, {80,160,320}, {80,160,400}, {90,180,360},
        {90,180,450}, {100,200,400}, {100,200,500}, {110,220,440}, {110,220,550},
        {120,240,480}, {120,240,600}, {130,260,520}, {130,260,650}, {140,280,560},
        {140,280,700}, {150,300,600}, {150,300,750}, {160,320,640}, {160,320,800}
    };

    std::vector<std::vector<int>> tiny = {
        {3, 5, 10}, {4, 6, 12}, {5, 8, 15}
    };  
    std::vector<std::vector<int>> baseline = {
        {10, 20, 40}, {15, 25, 50}, {20, 30, 60}, {25, 40, 80}
    };

    
    generate_group("tiny", tiny, rng);
    generate_group("baseline", baseline, rng);
    generate_group("small",  small,  rng);
    generate_group("medium", medium, rng);
    generate_group("large",  large,  rng);

    return 0;
}
\end{lstlisting}

\section{CPLEX implementacija}
Model je implementiran u C++ korišćenjem IBM ILOG CPLEX Concert Technology.
Vreme izvršavanja je ograničeno na 2h (7200s), a broj niti na 1 radi konzistentnosti.

\subsection*{Kod rešenja (CPLEX C++)}
\begin{lstlisting}[language=C++]
#include <fstream>
#include <ilcplex/ilocplex.h>
#include <iostream>
#include <iomanip>

ILOSTLBEGIN

int main(int argc, char* argv[]) {
    if (argc < 2) {
        std::cerr << "Usage: ./solve_tsfc instance.txt\n";
        return 1;
    }

    std::ifstream in(argv[1]);
    if (!in) {
        std::cerr << "Cannot open file.\n";
        return 1;
    }

    IloEnv env;

    IloInt I, J, K;
    in >> I >> J >> K;

    IloNumArray s(env, I);
    IloNumArray d(env, K);

    for (IloInt i = 0; i < I; i++) in >> s[i];
    for (IloInt k = 0; k < K; k++) in >> d[k];

    IloArray<IloNumArray> b(env, I);
    IloArray<IloNumArray> f(env, I);
    IloArray<IloNumArray> c(env, J);
    IloArray<IloNumArray> g(env, J);

    for (IloInt i = 0; i < I; i++) {
        b[i] = IloNumArray(env, J);
        f[i] = IloNumArray(env, J);
    }

    for (IloInt j = 0; j < J; j++) {
        c[j] = IloNumArray(env, K);
        g[j] = IloNumArray(env, K);
    }

    for (IloInt i = 0; i < I; i++)
        for (IloInt j = 0; j < J; j++) in >> b[i][j];

    for (IloInt i = 0; i < I; i++)
        for (IloInt j = 0; j < J; j++) in >> f[i][j];

    for (IloInt j = 0; j < J; j++)
        for (IloInt k = 0; k < K; k++) in >> c[j][k];

    for (IloInt j = 0; j < J; j++)
        for (IloInt k = 0; k < K; k++) in >> g[j][k];

    in.close();

    try {
        IloModel model(env);

        IloNumVarArray x(env, I * J, 0.0, IloInfinity, ILOFLOAT);
        IloBoolVarArray z(env, I * J);
        IloNumVarArray y(env, J * K, 0.0, IloInfinity, ILOFLOAT);
        IloBoolVarArray w(env, J * K);

        auto idx_x = [J](IloInt i, IloInt j) { return i * J + j; };
        auto idx_y = [K](IloInt j, IloInt k) { return j * K + k; };

        IloExpr obj(env);
        for (IloInt i = 0; i < I; i++)
            for (IloInt j = 0; j < J; j++)
                obj += b[i][j] * x[idx_x(i,j)] + f[i][j] * z[idx_x(i,j)];

        for (IloInt j = 0; j < J; j++)
            for (IloInt k = 0; k < K; k++)
                obj += c[j][k] * y[idx_y(j,k)] + g[j][k] * w[idx_y(j,k)];

        model.add(IloMinimize(env, obj));
        obj.end();

        for (IloInt i = 0; i < I; i++) {
            IloExpr expr(env);
            for (IloInt j = 0; j < J; j++)
                expr += x[idx_x(i,j)];
            model.add(expr - s[i] <= 0);
            expr.end();
        }

        for (IloInt k = 0; k < K; k++) {
            IloExpr expr(env);
            for (IloInt j = 0; j < J; j++)
                expr += y[idx_y(j,k)];
            model.add(expr - d[k] == 0);
            expr.end();
        }

        for (IloInt j = 0; j < J; j++) {
            IloExpr expr(env);
            for (IloInt i = 0; i < I; i++)
                expr += x[idx_x(i,j)];
            for (IloInt k = 0; k < K; k++)
                expr -= y[idx_y(j,k)];
            model.add(expr == 0);
            expr.end();
        }

        for (IloInt i = 0; i < I; i++)
            for (IloInt j = 0; j < J; j++)
                model.add(x[idx_x(i,j)] - s[i] * z[idx_x(i,j)] <= 0);

        for (IloInt j = 0; j < J; j++)
            for (IloInt k = 0; k < K; k++)
                model.add(y[idx_y(j,k)] - d[k] * w[idx_y(j,k)] <= 0);

        IloCplex cplex(model);
        cplex.setParam(IloCplex::Param::TimeLimit, 7200);
        cplex.setParam(IloCplex::Param::Threads, 1);
        cplex.setOut(env.getNullStream());

        if (cplex.solve() && cplex.isPrimalFeasible()) {
            std::cout << std::fixed << std::setprecision(2);
            std::cout << argv[1] << " "
                      << ((cplex.getStatus() == IloAlgorithm::Optimal) ? "OPTIMAL" : "FEASIBLE") << " "
                      << cplex.getObjValue() << " "
                      << cplex.getTime() << " "
                      << cplex.getNiterations() << " "
                      << cplex.getNnodes() << "\n";
        }

    } catch (IloException &e) {
        std::cerr << "CPLEX Exception: " << e << std::endl;
    }

    env.end();
    return 0;
}
\end{lstlisting}

\section{Rezultati CPLEX-a}

Vreme izvršavanja je ograničeno na 2 sata po instanci.
U tabelama su prikazani:
\begin{itemize}
    \item status (OPTIMAL ili FEASIBLE),
    \item vrednost ciljne funkcije,
    \item vreme izvršavanja (s),
    \item broj iteracija,
    \item broj čvorova.
\end{itemize}

\subsection{Male instance (small)}
\begin{longtable}{l l r r r r}
\toprule
\textbf{Instanca} & \textbf{Status} & \textbf{Obj} & \textbf{Vreme (s)} & \textbf{Iteracije} & \textbf{Čvorovi} \\
\midrule
\endhead

small\_1  & OPTIMAL  & 150568.00 & 1.30   & 494      & 13 \\
small\_2  & OPTIMAL  & 156022.00 & 1.51   & 969      & 48 \\
small\_3  & OPTIMAL  & 161224.00 & 1.92   & 745      & 60 \\
small\_4  & OPTIMAL  & 186537.00 & 1.90   & 467      & 11 \\
small\_5  & OPTIMAL  & 213785.00 & 2.06   & 754      & 18 \\
small\_6  & OPTIMAL  & 188954.00 & 7.15   & 9348     & 583 \\
small\_7  & OPTIMAL  & 231984.00 & 17.97  & 22618    & 1308 \\
small\_8  & OPTIMAL  & 202467.00 & 11.79  & 8786     & 367 \\
small\_9  & OPTIMAL  & 210740.00 & 25.79  & 28303    & 1551 \\
small\_10 & OPTIMAL  & 258236.00 & 7.38   & 124646   & 5628 \\
small\_11 & OPTIMAL  & 286991.00 & 84.08  & 1168471  & 37922 \\
small\_12 & FEASIBLE & 284643.00 & 1818.90& 5118195  & 119052 \\
small\_13 & OPTIMAL  & 278044.00 & 43.03  & 100035   & 4193 \\
small\_14 & OPTIMAL  & 284648.00 & 2509.23& 7374762  & 207598 \\
small\_15 & OPTIMAL  & 315417.00 & 1585.38& 3325528  & 71013 \\
small\_16 & FEASIBLE & 349820.00 & 3479.76& 6251954  & 101284 \\
small\_17 & FEASIBLE & 368942.00 & 2606.42& 5470859  & 128331 \\
small\_18 & FEASIBLE & 349036.00 & 4359.04& 8837101  & 171084 \\
small\_19 & FEASIBLE & 351345.00 & 4018.04& 7789292  & 142852 \\
small\_20 & FEASIBLE & 400600.00 & 3203.13& 5178550  & 110993 \\
\bottomrule
\end{longtable}

\subsection{Srednje instance (medium)}
U nastavku su prikazane instance za koje su rezultati dostupni:

\begin{center}
\begin{tabular}{l l r r r r}
\toprule
\textbf{Instanca} & \textbf{Status} & \textbf{Obj} & \textbf{Vreme (s)} & \textbf{Iteracije} & \textbf{Čvorovi} \\
\midrule
medium\_11 & FEASIBLE & 551427.00 & 4242.68 & 8148431 & 228474 \\
medium\_12 & FEASIBLE & 478570.00 & 7194.29 & 9805905 & 250345 \\
\bottomrule
\end{tabular}
\end{center}

\subsection{Baseline instance}
U eksperimentu su testirane i baseline instance:

\begin{center}
\begin{tabular}{l l r r r r}
\toprule
\textbf{Instanca} & \textbf{Status} & \textbf{Obj} & \textbf{Vreme (s)} & \textbf{Iteracije} & \textbf{Čvorovi} \\
\midrule
baseline\_1 & OPTIMAL  & 223136.00 & 1.70    & 14281     & 404 \\
baseline\_2 & OPTIMAL  & 317975.00 & 1190.98 & 11745067  & 265058 \\
baseline\_3 & OPTIMAL  & 322093.00 & 3695.07 & 21479433  & 385399 \\
baseline\_4 & FEASIBLE & 399062.00 & 7189.11 & 1280093   & 31964 \\
\bottomrule
\end{tabular}
\end{center}

\section{Komentar i analiza rezultata}
Na osnovu dobijenih rezultata može se uočiti sledeće:
\begin{itemize}
    \item Za manje instance CPLEX u velikom broju slučajeva pronalazi optimalno rešenje relativno brzo, ali kako dimenzije rastu, broj čvorova i iteracija značajno raste.
    
    \item Kod nekih instanci (npr. \texttt{small\_12}, \texttt{small\_16}-\texttt{small\_20}) CPLEX nije uspeo da dokaže optimalnost u okviru vremenskog limita, ali je pronašao dopustivo rešenje (status FEASIBLE).
    
    \item Za srednje instance prikazane u tabeli (\texttt{medium\_11 i medium\_12}), CPLEX je dostigao vremenski limit i dao samo dopustiva rešenja, što ukazuje da problem postaje veoma težak već u ovoj dimenziji.
    
    \item Vreme izvršavanja je u korelaciji sa rastom dimenzija i kompleksnošću MILP modela zbog binarnih promenljivih $z_{ij}$ i $w_{jk}$, koje uvode kombinatornu eksploziju u pretrazi.

\end{itemize}

\begin{tcolorbox}[myrednote,title=Napomena o rezultatima većih instanci]
U trenutnoj verziji rada nisu prikazani rezultati za sve instance iz skupa \texttt{medium}, kao nijedna iz skupa \texttt{large}. 
Tokom izvršavanja eksperimenta pojavili su se tehnički problemi prilikom rešavanja pojedinih instanci: na više različitih mašina dolazilo je do prekida rada računara, dok je za instancu \texttt{medium\_10} program u više pokušaja neočekivano prekidao izvršavanje bez ispisivanja rezultata. 
Zbog navedenog, u radu su trenutno prikazani samo rezultati za instance za koje je izvršavanje uspešno završeno. 
Planirano je da se preostali eksperimenti ponove na drugoj konfiguraciji, nakon čega će rezultati biti naknadno dopunjeni.
\end{tcolorbox}

\begin{tcolorbox}[myrednote,title=Napomena o rezultatima i pribavljanju referentnih instanci]
U okviru ovog rada nije bilo moguće izvršiti poređenje dobijenih rezultata sa referentnim vrednostima iz literature, jer originalne instance korišćene u radu autora nisu bile dostupne za preuzimanje. Iako je u literaturi naveden link ka instancama, nakon višenedeljnih pokušaja pristupa i slanja upita, kao i direktnog kontaktiranja istraživača koji su sproveli istraživanje, instance nisu pribavljene. 
\end{tcolorbox}


\section{Link za kod, izvrsnu verziju i korišćene instance}
Kod programa i izvršna verzija (CPLEX) dostupni su na sledećem linku:

\begin{center}
\url{https://github.com/MSZ2/TS-FCTP-Matheuristics-project/tree/CPLEX}
\end{center}

\begin{thebibliography}{9}

\bibitem{tsfctp}
H.~I.~Calvete, C.~Gal\'e, J.~A.~Iranzo, P.~Toth,
\textit{A matheuristic for the two-stage fixed-charge transportation problem}.

\end{thebibliography}


\end{document}
